\section{Conclusiones}

Este estudio examina el diseño de las metas de gestión en la regulación peruana de los servicios de agua potable y saneamiento. El problema central es que el Índice de Cumplimiento Global (ICG) puede combinar indicadores que miden objetos regulatorios distintos: resultados observables de la prestación, variables de eficiencia o gestión, obligaciones de implementación y ejecución financiera de inversiones o reservas. Esta composición no es neutral, pues condiciona tanto la interpretación del índice como los incentivos que transmite a la empresa prestadora.

La revisión del marco regulatorio y la evidencia de las fiscalizaciones realizadas por la Sunass durante 2024 muestran una marcada heterogeneidad entre las metas evaluadas. El cumplimiento promedio general asciende a 71\%, mientras que los menores niveles se concentran principalmente en metas vinculadas con inversiones, reservas y otras obligaciones de ejecución. Esta evidencia es descriptiva y no identifica causalmente las razones de las diferencias observadas, pero muestra que un mismo porcentaje de incumplimiento puede contener información distinta según la naturaleza del objeto medido, su horizonte temporal y los factores que condicionan su ejecución.

A partir de esta heterogeneidad, el análisis distingue entre los resultados de desempeño comprendidos en \(J^{Q}\), las inversiones vinculadas con su producción, \(R^{Q}\), las demás obligaciones discretas de implementación, \(R^{O}\), y los indicadores de gestión agrupados en \(C^{O}\), cuya tecnología de producción no se representa mediante una inversión específica. Esta clasificación permite separar analíticamente la medición de los resultados, la verificación de las obligaciones de implementación y el seguimiento de otras dimensiones regulatorias sin restar relevancia ni exigibilidad a ninguna de ellas.

La comparación secuencial \(A_{0}\rightarrow A\rightarrow A'\rightarrow A''\rightarrow B\) muestra que los problemas identificados corresponden a mecanismos distintos. La transición \(A_{0}\rightarrow A\) focaliza la medición financiera y elimina, para las inversiones de \(R^{Q}\), la superposición entre su contribución a los resultados y su ejecución financiera dentro del índice, aunque ello no garantiza por sí mismo un mayor esfuerzo contemporáneo de ejecución. La transición \(A\rightarrow A'\) sustituye la ejecución financiera por el alcance físico o funcional como base de verificación de \(R^{O}\) y evita que una inversión completada por debajo del presupuesto continúe siendo tratada como incumplida por el solo hecho de generar un ahorro.

La transición \(A'\rightarrow A''\) corrige un mecanismo diferente. Cuando la consecuencia asociada con las inversiones depende de un umbral agregado del ICG, el cumplimiento de unas dimensiones puede desactivar la sanción aun cuando subsistan obligaciones pendientes. Las sanciones independientes preservan, en cambio, la consecuencia asociada con cada inversión mientras permanezca incumplida, sin perjuicio de la evaluación de su atribuibilidad y de las demás condiciones exigidas por el régimen sancionador.

Finalmente, la transición \(A''\rightarrow B\) concentra el ICG en los resultados de \(J^{Q}\). La exclusión de la ejecución financiera y de los indicadores agrupados en \(C^{O}\) evita que componentes que representan otros objetos regulatorios reduzcan el peso efectivo de los resultados de cobertura y calidad dentro del índice. Esta delimitación no implica que la ejecución de inversiones, la eficiencia empresarial, la sostenibilidad financiera u otras dimensiones de gestión dejen de ser observadas, sino que su seguimiento se realiza mediante instrumentos diferenciados y conserva una interpretación separada.

Los resultados permiten también precisar la relación entre fiscalización y sanción. Cuando el objeto regulatorio es asegurar la ejecución de una inversión, la verificación del cumplimiento debería distinguir entre alcance físico o funcional, costo y entrada efectiva en operación. Una inversión completada con un costo inferior al presupuesto no constituye necesariamente un incumplimiento si conserva el alcance, la calidad, la capacidad y las demás condiciones exigidas. De manera complementaria, la base económica de una eventual sanción debería guardar correspondencia con la obligación efectivamente incumplida y evitar que una misma inversión sea valorizada repetidamente cuando contribuya a varias metas o resultados.

Los incumplimientos de resultados requieren una lógica distinta. Una condición insuficiente de continuidad, presión, acceso, calidad sanitaria u otra dimensión de \(J^{Q}\) puede justificar medidas correctivas, compensaciones o tipificaciones específicas cuando así lo establezca la regulación, pero su consecuencia no tiene por qué utilizar automáticamente como base económica el costo de una inversión relacionada. Del mismo modo, los indicadores comprendidos en \(C^{O}\) pueden conservar metas, supervisión y consecuencias autónomas cuando corresponda, sin necesidad de incorporarlos al ICG de desempeño.

La extensión \(B^{+}\) examina una cuestión adicional: el tratamiento del sobrecumplimiento. Mantener el ICG ordinario truncado en 100\% evita que una mejora extraordinaria en una dimensión compense el incumplimiento de otra, mientras que un mecanismo separado puede preservar un incentivo marginal para producir resultados adicionales. Su conveniencia, sin embargo, requiere condiciones más exigentes de verificabilidad, controlabilidad, sostenibilidad y no desplazamiento de esfuerzos desde obligaciones pendientes, por lo que no constituye una consecuencia necesaria de la separación del ICG.

Los resultados no establecen una superioridad incondicional de una arquitectura sobre todas las demás. La conveniencia de separar instrumentos depende de que la fiscalización y el régimen sancionador preserven incentivos suficientes para cumplir las obligaciones de implementación, de que los indicadores incorporados al ICG representen adecuadamente los resultados que se pretende hacer visibles y de que la señal de desempeño alcance suficiente difusión y legibilidad. La comparación de bienestar considera, además, el gasto efectivo de inversión, los costos de operación y mantenimiento y la carga sectorial de las multas, que puede diferir de su valor nominal cuando la sanción constituye parcialmente una transferencia. Asimismo, el modelo es una representación estilizada de los incentivos y no identifica las respuestas dinámicas que podrían producirse ante una reforma.

En conjunto, el análisis sugiere que la pertenencia de distintas obligaciones a un mismo régimen de metas de gestión no exige que todas incidan de la misma manera en una medida agregada. Diferenciar entre resultados de desempeño, obligaciones de implementación, indicadores de gestión y mejoras adicionales permite preservar la identidad del objeto regulado, reducir la superposición entre instrumentos y hacer más interpretable la señal proporcionada por el ICG. La principal implicancia no es sustituir un instrumento por otro, sino asignar a cada función regulatoria un mecanismo congruente con la conducta o el resultado que se pretende medir, supervisar o disciplinar.